% Part: sets-functions-relations
% Chapter: relations
% Section: special-properties

\documentclass[../../../include/open-logic-section]{subfiles}

\begin{document}

\olfileid[fa-IR]{sfr}{rel}{prp}
\olsection{ویژگی‌های خاص روابط}

\begin{intro}
برخی گونه‌های روابط چنان رایج‌اند که نام‌هایی خاص به آنها داده‌اند.
برای نمونه، $\le$ و~$\subseteq$ هر دو اعضای دامنه‌های مربوط به خود را
(مثلاً $\Nat$ در مورد~$\le$ و $\Pow{A}$ در مورد~$\subseteq$) به
شیوه‌هایی مشابه به هم مرتبط می‌کنند. برای آنکه دقیقاً دریابیم این روابط
از چه جهت مشابه و از چه جهت متفاوت‌اند، آنها را بر پایهٔ برخی ویژگی‌های
خاصی که روابط می‌توانند داشته باشند دسته‌بندی می‌کنیم. معلوم می‌شود که
(ترکیب‌هایی از) بعضی از این ویژگی‌های خاص اهمیت ویژه‌ای دارند: ترتیب‌ها
و روابط هم‌ارزی.
\end{intro}

\begin{defn}[بازتابی‌بودن]
رابطهٔ $R \subseteq A^2$ \emph{بازتابی} است اگر و تنها اگر برای هر $x \in
A$، $Rxx$.
\end{defn}

\begin{defn}[تراگذری]
رابطهٔ $R \subseteq A^2$ \emph{تراگذری} است اگر و تنها اگر هرگاه $Rxy$
و $Ryz$، آنگاه $Rxz$ نیز.
\end{defn}

\begin{defn}[تقارن]
رابطهٔ~$R \subseteq A^2$ \emph{متقارن} است اگر و تنها اگر هرگاه
$Rxy$، آنگاه~$Ryx$ نیز.
\end{defn}

\begin{defn}[پادتقارنی]
رابطهٔ~$R \subseteq A^2$ \emph{پاد\-مت\-قارن} است اگر و تنها اگر هرگاه هر دو
$Rxy$ و $Ryx$ برقرار باشند، آنگاه $x=y$ (یا به بیان دیگر: اگر $x\neq y$،
آنگاه دست‌کم یکی از $\lnot Rxy$ یا $\lnot Ryx$ برقرار است).
\end{defn}

\begin{explain}
در رابطه‌ای متقارن، $Rxy$ و $Ryx$ همیشه با هم برقرارند یا هیچ‌یک برقرار
نیست. در رابطه‌ای پادمتقارن، تنها در صورتی ممکن است $Rxy$ و $Ryx$ با هم
برقرار باشند که $x = y$. توجه کنید که این امر \emph{ایجاب نمی‌کند} که
$Rxy$ و $Ryx$ وقتی $x = y$ است برقرار باشند؛ فقط برقرار بودنشان را منتفی
نمی‌کند. پس یک رابطهٔ پادمتقارن می‌تواند بازتابی باشد، اما هر رابطهٔ
پادمتقارنی بازتابی نیست. همچنین توجه کنید که پادمتقارن‌بودن با صرفاً
متقارن‌نبودن دو شرط متفاوت‌اند. در واقع، یک رابطه می‌تواند هم‌زمان هم
متقارن و هم پادمتقارن باشد (برای نمونه، رابطهٔ همانی چنین است).
\end{explain}

\begin{defn}[همبندی]
رابطهٔ $R \subseteq A^2$ \emph{همبند} است اگر برای همهٔ $x,y\in
A$، اگر $x \neq y$، آنگاه دست‌کم یکی از $Rxy$ و~$Ryx$ برقرار است.
\end{defn}

\begin{prob}
نمونه‌هایی از روابطی به دست دهید که (a) بازتابی و متقارن ولی
تراگذری نیستند، (b) بازتابی و پادمتقارن‌اند، (c) پادمتقارن و تراگذری‌اند
ولی بازتابی نیستند، و (d) بازتابی، متقارن و تراگذری‌اند. از روابط روی اعداد
یا مجموعه‌ها استفاده نکنید.
\end{prob}

\begin{defn}[غیربازتابی‌بودن]
رابطهٔ $R \subseteq A^2$ \emph{غیربازتابی} نامیده می‌شود اگر برای هر $x \in
A$، نه $Rxx$.
\end{defn}

\begin{defn}[نامتقارنی]
رابطهٔ $R \subseteq A^2$ \emph{نامتقارن} نامیده می‌شود اگر هیچ زوج $x,y\in
A$ وجود نداشته باشد که هر دو $Rxy$ و~$Ryx$ برقرار باشند.
\end{defn}

توجه کنید که اگر $A \neq \emptyset$، آنگاه هیچ رابطهٔ غیربازتابی
روی~$A$ بازتابی نیست و هر رابطهٔ نامتقارن روی~$A$ پادمتقارن نیز هست.
بااین‌حال، روابطی $R \subseteq A^2$ وجود دارند که نه بازتابی‌اند و نه
غیربازتابی، و روابط پادمتقارنی نیز وجود دارند که نامتقارن نیستند.

\end{document}
